\section{Introduction}
\label{sec:intro}

Budgeted test-time inference is often summarized by one performance number: final accuracy under a
shared nominal budget of samples, branches, or reasoning actions. For performance evaluation, that
summary is incomplete because the nominal budget is an experimental entitlement, not an observed
resource measurement. Accuracy under the same entitlement can therefore conflate several distinct
questions:
\begin{itemize}
\item Did any correct answer enter the candidate pool?
\item Did the commit rule select a correct answer when one was present?
\item Did compared methods share an identical pool, or did one simply receive more diverse
candidates?
\item Does a failure-derived selection rule that helps under one commercial API transfer to another?
\item Does a matched nominal budget $B$ imply comparable calls, tokens, latency, or cost?
\end{itemize}

The distinction matters because equal nominal budgets need not induce equal realized resources.
A method may spend the same logical budget while producing different numbers of successful
completions, retries, generated tokens, observed latency, or estimated cost. When those quantities
are collapsed into final accuracy, a selector improvement can be confused with better candidate
discovery, and a budget-matched comparison can be mistaken for a resource-matched one.

Closed commercial APIs make the measurement problem harder. Providers expose different operational
behavior, telemetry fields may be incomplete, and hosted model identifiers do not guarantee
immutable weight-level versions. A post-generation rule can also look effective when compared only
to a single generator, then lose to plurality on the \emph{same} extracted answers. Failed or
unscorable provider--dataset cells create another reporting choice: they can be imputed, dropped, or
preserved as protocol outcomes.

Prior work contains several relevant ingredients individually: coverage and pass@$k$ analyses,
fixed-pool diagnostics, and self-consistency or other consensus aggregators
\citep{passatk,ll_monkeys,oracle_gap,selfconsistency}. Router and test-time-compute work also
motivates resource reporting under budgeted inference \citep{snell_ttc}. The measurement gap is not
the absence of any one ingredient. It is the lack of one executable and auditable closed-API
protocol that forces these separations to be reported together for the same provider--dataset
matrix.

\begin{figure}[t]
\centering
\resizebox{\linewidth}{!}{\begin{tikzpicture}[
  font=\small,
  node distance=0.9cm and 0.55cm,
  box/.style={draw, rounded corners=1.5pt, align=center, inner sep=5pt, minimum height=1.05cm},
  wide/.style={box, text width=0.19\linewidth},
  smallbox/.style={box, text width=0.17\linewidth},
  arrow/.style={-{Latex[length=2mm]}, line width=0.45pt}
]
\node[wide] (instance) {Benchmark\\instance};
\node[wide, right=of instance] (generation) {Class A generation\\nominal budget $B$};
\node[wide, right=of generation] (pool) {Extracted candidate\\pool};
\node[wide, right=of pool] (selection) {Class B selection\\same pool, $0$ extra calls};

\draw[arrow] (instance) -- (generation);
\draw[arrow] (generation) -- (pool);
\draw[arrow] (pool) -- (selection);

\node[smallbox, below=of generation] (resources) {Realized resource\\vector $\mathbf{r}$:\\calls, tokens,\\latency, cost};
\node[smallbox, below=of pool] (discovery) {Discovery\\correct answer present?};
\node[smallbox, below=of selection] (commit) {Commit outcome\\selected answer correct?};

\draw[arrow] (generation) -- (resources);
\draw[arrow] (pool) -- (discovery);
\draw[arrow] (selection) -- (commit);

\node[box, text width=0.31\linewidth, below=1.0cm of discovery] (cells) {Provider $\times$ dataset cells:\\completed, transfer-labeled, or protocol-blocked};
\draw[arrow] (resources.south) |- (cells.west);
\draw[arrow] (commit.south) |- (cells.east);
\end{tikzpicture}
}
\caption{Protocol schematic. Nominal budget $B$ defines the logical action contract; realized
resources are measured separately. Selection claims are attributed after candidate availability is
fixed, and provider$\times$dataset cells may be completed, transfer-labeled, or protocol-blocked.}
\label{fig:protocol}
\end{figure}

\paragraph{Contributions}
This article contributes a controlled performance-evaluation protocol for budgeted closed-API
inference (Figure~\ref{fig:protocol}). The contribution is the joint accounting contract: the same
evaluation records nominal entitlements, realized resources, discovery, selector attribution,
provider-transfer labels, provenance tiers, protocol-blocked cells, and telemetry completeness, so
each performance conclusion has an explicit measurement layer. Specifically:
\begin{itemize}
\item \textbf{Protocol formulation.} We define nominal budget, candidate-pool discovery,
conditional selection accuracy, consensus regret, realized resource vectors, and protocol-blocked
outcomes for closed-API inference.
\item \textbf{Selector attribution and transfer analysis.} We compare nontrivial commit rules to an
identical-pool plurality control and label failure-mined rules by development and transfer
environment.
\item \textbf{Realized-resource accounting.} We separate successful completions, retries, tokens,
latency, and estimated dollars, while marking incomplete telemetry and reconstruction status.
\item \textbf{Released evidence boundary.} We report the blocked cell rather than imputing it and
provide aggregate-record verification artifacts for the reported tables and numerical invariants.
\end{itemize}

The empirical study exercises the protocol on four providers and four reasoning workloads at one
nominal budget, with 15 completed cells and one protocol-blocked cell. It demonstrates that
identical-pool plurality can change selector conclusions, that a failure-mined gate can transfer
poorly across providers, and that reconstructed successful completions can change efficiency
interpretation even when accuracy and paired tests are unchanged.

The study does not claim a general algorithm ranking, a complete cost-efficiency or Pareto analysis,
official reproductions of L1/S1/TALE, or guaranteed regeneration of identical proprietary API
outputs. Frontier, FTA, Pooled-4, and the adapters are measurement objects used to exercise the
protocol, not the central contribution.

\paragraph{Roadmap}
Section~\ref{sec:related} positions the work relative to performance measurement, test-time
compute, routing, consensus, and verifier literature. Sections~\ref{sec:problem}--\ref{sec:setup}
define the evaluation objects and closed-API workload. Sections~\ref{sec:results}--\ref{sec:compute}
report the matrix, identical-pool selection tests, provider transfer, and resource accounting.
Section~\ref{sec:ablations} reports offline counterfactuals, and
Sections~\ref{sec:limitations}--\ref{sec:conclusion} summarize scope and implications for
performance-evaluation practice.

\paragraph{Measurement objects}
The named procedures are measurement objects for exercising the protocol
(Section~\ref{sec:method}). Class~A single-generator policies under $B{=}6$ include a
reserved-attempt full-expansion diagnostic (Frontier) and closed-API adapters inspired by L1, S1,
and TALE \citep{l1,s1,tale}. Class~B post-generation selectors over the shared pool include
Pooled-4 and FTA. Offline oracles are upper bounds only. A single-cell Azure$\times$GSM8K
call/sample-budget control compares Frontier to same-model self-consistency with $N{=}6$
(SC-$N{=}6$; Section~\ref{sec:sc-entitlement-control}).
